% =============================================================
% TDSC published-look starter (visual style close to IEEE Xplore)
% NOTE: IEEE may apply its own production typesetting after acceptance.
% For final submission, also compare with the current IEEE Template Selector.
% =============================================================
\documentclass[10pt,journal]{IEEEtran}

\usepackage{amsmath,amssymb}
\usepackage{graphicx}
\usepackage{booktabs}
\usepackage{multirow}
\usepackage{cite}
\usepackage[caption=false,font=footnotesize]{subfig}
\usepackage{url}

\begin{document}

\title{Active Diagnostic Probing for Backdoor Sample Cleansing in Contrastive Learning}

\author{First~Author,~\IEEEmembership{Graduate Student Member,~IEEE,}
        Second~Author,~\IEEEmembership{Member,~IEEE,}
        and~Third~Author,~\IEEEmembership{Senior Member,~IEEE}%
\thanks{First Author and Second Author are with the School of XXX, XXX University, City, Country. E-mail: xxx@xxx.edu.}%
\thanks{Third Author is with the Department of XXX, XXX University, City, Country.}%
\thanks{Manuscript received Month xx, 2026; revised Month xx, 2026.}}


\maketitle

\begin{abstract}
Contrastive learning has emerged as an effective paradigm for learning transferable representations from large-scale unlabeled data. However, its reliance on uncurated pre-training data exposes it to data-poisoning backdoor attacks. In this work, we study backdoored sample cleansing for unlabeled contrastive-learning datasets and develop an active diagnostic probing framework. The key idea is to transform difficult-to-separate static representations into probe-induced response representations, calibrate benign responses using a small cross-domain reference set, preserve local response topology with UMAP, and identify suspicious samples with reference-anchored DBSCAN. Extensive experiments will evaluate cleansing effectiveness, post-cleansing security and utility, and computational overhead.
\end{abstract}

\begin{IEEEkeywords}
Contrastive learning, self-supervised learning, backdoor attacks, poisoned sample detection, data cleansing.
\end{IEEEkeywords}

\section{Introduction}\label{sec:introduction}
\IEEEPARstart{C}{ontrastive} learning (CL) has become a central paradigm for learning transferable representations from large-scale unlabeled data. Representative methods learn invariance across augmented views through instance discrimination, momentum-encoded dictionaries, or bootstrapped prediction~\cite{Chen2020SimCLR,He2020MoCo,Grill2020BYOL}. By reducing the dependence on manual annotation, CL enables pre-trained encoders to support a broad range of downstream tasks. This advantage, however, also expands the attack surface: pre-training data are often collected from large, weakly curated sources in which an adversary can inject a small number of poisoned samples. The resulting encoder may behave normally on benign inputs yet map trigger-bearing inputs toward an attacker-selected semantic region, causing downstream models to inherit a hidden backdoor~\cite{Carlini2022Poisoning,Saha2022SSLBackdoor,Liu2022PoisonedEncoder}. Furthermore, the benign and backdoor objectives are intertwined during contrastive training, yielding learning dynamics and feature distributions that differ from those of supervised backdoor attacks; consequently, defenses transferred directly from supervised learning can be ineffective~\cite{Li2023CTRL,Li2024Difficulty}.

This work considers a preventive defense at the data level: identifying and removing backdoored samples from an unlabeled pre-training corpus before a final encoder is trained from scratch. This objective differs from detecting or repairing a released encoder, as pursued by DECREE and SSL-Cleanse~\cite{Feng2023DECREE,Zheng2024SSLCleanse}. Existing data-oriented defenses reveal several useful directions. PatchSearch exploits the localized influence of patch triggers and trains a poison classifier from highly suspicious regions~\cite{Tejankar2023PatchSearch}. Confusion Training and ASSET actively enlarge behavioral differences between benign and poisoned samples through additional training or optimization~\cite{Qi2023Proactive,Pan2023ASSET}. AUTO bootstraps relatively benign and relatively poisonous subsets from the contaminated corpus and iteratively refines their supervisory signals~\cite{Chen2025AUTO}. These studies demonstrate the feasibility of sample cleansing, but also expose a tension among trigger generality, feature separability, optimization cost, and the assumptions placed on trusted data.

A central difficulty is that poisoned samples need not be isolated in the original representation space. They retain genuine semantic content, while the malicious trigger--target association is learned jointly with benign invariances. Recent attacks establish this association using substantially different mechanisms, including visible patches, frequency-domain perturbations, augmentation-aware poison construction, and bi-level trigger optimization~\cite{Saha2022SSLBackdoor,Li2023CTRL,Zhang2024CorruptEncoder,Sun2024BLTO}. Static embeddings or passive clustering statistics can therefore overlap across benign and poisoned samples, and a signature tailored to one trigger family may not transfer to another~\cite{Pan2023ASSET,Li2024Difficulty}. Patch localization is effective when a compact spatial trigger dominates the response, but its applicability narrows for distributed or optimized perturbations~\cite{Tejankar2023PatchSearch}. Active optimization can improve separation, yet it may require repeated training and reliable supervisory signals~\cite{Qi2023Proactive,Pan2023ASSET,Chen2025AUTO}. These limitations motivate an intervention that exposes how a sample reacts to a controlled semantic perturbation instead of assuming that its unperturbed embedding is already anomalous.

Our key hypothesis is that the latent backdoor association constrains a poisoned sample's response to a second, controlled semantic signal. We repurpose a pre-trained BLTO trigger generator as a diagnostic probe~\cite{Sun2024BLTO}. The probe is associated with a defender-selected reference semantic, whereas the attacker's target is chosen independently and may be either different from or identical to that reference. When the two targets differ, the diagnostic probe and the original backdoor exert competing semantic attraction; benign and poisoned samples are therefore expected to exhibit different probe-conditioned geometry. This is a testable mechanism hypothesis rather than a secrecy assumption: the defense does not rely on hiding the probe generator or its reference semantic, and the same-target case remains an explicit boundary condition to be evaluated.

To make the reference signal available during representation learning, the defender augments the contaminated corpus with a small trusted reference set after applying the diagnostic probe to those reference images. The probed references and the unlabeled training samples undergo the same augmentations and participate in the same unsupervised contrastive objective; reference labels are used only to select one semantic category and never enter a supervised loss. We set the reference budget to approximately $1\%$ of the training corpus and reuse one pre-trained probe generator across datasets and encoders. After this diagnostic training stage, the original and probe-transformed views of each training sample provide complementary representations, while the probed trusted samples identify the benign reference region without requiring their source distribution to reproduce the entire unlabeled corpus.

The resulting framework jointly embeds the training and reference representations with UMAP to preserve local neighborhood structure~\cite{McInnes2018UMAP}, and then applies DBSCAN in the low-dimensional space~\cite{Ester1996DBSCAN}. The mean location of all reference points anchors the decision: the non-noise cluster whose training-sample center is closest to the reference center is retained as benign, whereas disconnected clusters and noise are treated as suspicious. The cleaned corpus is subsequently used to train the final encoder from scratch without the reference set. This pipeline transforms a controlled probe response into a data-cleansing decision while avoiding trigger localization, poison labels, a known poisoning ratio, or repeated training of a new probe for every dataset.

The main contributions of this work are summarized as follows:
\begin{itemize}
    \item We formulate unlabeled backdoor sample cleansing as an active diagnostic problem and repurpose a reusable BLTO generator as a controlled semantic probe for exposing backdoor-constrained behavior.
    \item We incorporate a small, single-semantic trusted reference set into unsupervised diagnostic training, without using its labels as supervision or requiring it to match the full training distribution.
    \item We combine original and probe-conditioned representations with joint UMAP embedding and a reference-anchored DBSCAN decision rule to remove disconnected clusters and noise without knowing the poisoning ratio.
    \item We define an end-to-end evaluation protocol covering detection quality, post-cleansing security and utility, reference-set size, probe--attack target alignment, computational overhead, and stochastic stability.
\end{itemize}

\section{Related Work}\label{sec:relatedwork}
\subsection{Backdoor Attacks on Contrastive Learning}
Backdoor attacks on contrastive learning exploit the fact that representations, rather than task-specific decision boundaries, are learned during unlabeled pre-training. The adversary poisons a small subset of the pre-training corpus so that the encoder preserves normal utility on benign inputs while associating a trigger with an attacker-selected semantic target. Because downstream classifiers reuse the compromised representation, the association can survive task-specific fine-tuning or linear evaluation. Early studies established that contrastive objectives are vulnerable to both targeted poisoning and persistent backdoor behavior~\cite{Carlini2022Poisoning}. SSL-Backdoor stamps a patch on images from the target category, causing a model trained without labels to transfer the trigger--target association to downstream classification~\cite{Saha2022SSLBackdoor}. PoisonedEncoder further formulates poison construction as an optimization problem and supports multiple target inputs or target classes using only poisoned pre-training samples~\cite{Liu2022PoisonedEncoder}.

Subsequent work broadens both the trigger space and the attack mechanism. CTRL embeds a trigger in the frequency domain and exploits the invariance encouraged by contrastive training, demonstrating effective attacks at low poisoning ratios~\cite{Li2023CTRL}. CorruptEncoder analyzes how random cropping and positive-pair construction affect data-poisoning attacks, then designs poisoned samples that remain effective under contrastive augmentations~\cite{Zhang2024CorruptEncoder}. BLTO jointly optimizes a trigger and a surrogate contrastive-learning process: its inner problem simulates representation learning, whereas its outer problem maintains the association between triggered samples and the target semantics~\cite{Sun2024BLTO}. These attacks differ in visibility, spatial locality, and dependence on data augmentation, so the presence of a backdoor cannot be reduced to a single patch or frequency signature.

More fundamentally, the benign and malicious tasks are coupled differently in contrastive and supervised learning. The analysis in~\cite{Li2024Difficulty} shows that contrastive representations can entangle the backdoor objective with useful semantics, making a poisoned sample neither a conventional input outlier nor an isolated feature cluster. This observation explains why a defense that relies only on static embeddings, loss values, or a supervised classifier's decision statistics may transfer poorly across contrastive algorithms and attacks. Our work therefore treats the original embedding as an insufficient diagnostic view and studies how samples respond to a controlled, attack-informed probe.

\subsection{Backdoor Defenses for Contrastive Learning}
Existing defenses can be divided into model-oriented inspection and data-oriented cleansing. Model-oriented methods determine whether a released encoder is compromised or mitigate the backdoor after training. DECREE searches for a small perturbation that makes diverse inputs share similar representations and uses the recovered perturbation to detect backdoored encoders~\cite{Feng2023DECREE}. SSL-Cleanse combines representation clustering, trigger inversion, and unlearning to detect and mitigate Trojans in self-supervised models~\cite{Zheng2024SSLCleanse}. These methods are valuable when the encoder is the only available artifact, but their primary output is a model-level diagnosis or repaired encoder rather than a sanitized version of the original pre-training corpus.

Data-oriented defenses instead identify suspicious training samples so that a new encoder can be trained on a cleaner dataset. PatchSearch trains an initial encoder on the contaminated corpus, localizes salient candidate patches, evaluates their ability to induce anomalous behavior when pasted onto other images, and iteratively trains a poison classifier~\cite{Tejankar2023PatchSearch}. Its localization signal is well matched to compact patch triggers but is less directly applicable to distributed, frequency-domain, or sample-specific patterns. Confusion Training actively introduces additional label noise to decouple benign and backdoor correlations and thereby expose poisoned samples~\cite{Qi2023Proactive}. Although developed in a supervised setting, it motivates the broader principle that a diagnostic intervention can reveal a latent backdoor more clearly than passive observation.

ASSET extends active separation across supervised, self-supervised, and transfer-learning paradigms by optimizing an offset between the untrusted data and a small clean base set~\cite{Pan2023ASSET}. Its trusted anchor improves identifiability, but distribution mismatch between the base and training sets can become a competing separation signal. AUTO removes the external clean-data requirement by mining relatively benign and relatively poisonous subsets from the contaminated corpus and bootstrapping their refinement~\cite{Chen2025AUTO}. This reduces reliance on curated clean data, while making the result dependent on the quality of the initially mined supervision and iterative optimization. Our framework follows the active-separation principle but assigns the trusted data a different role. A small, single-semantic reference set is first transformed by a reusable diagnostic probe and participates in the same unsupervised contrastive training as the untrusted corpus. Its probe-conditioned representations subsequently anchor the benign cluster in the joint embedding. Thus, the reference set neither supplies poison labels nor attempts to reproduce the complete distribution of the training data.

\subsection{Representation-Based Anomaly Detection}
Representation-space analysis is a major line of backdoor sample detection. Spectral Signatures centers deep features and removes examples with large projections onto a dominant singular direction associated with poisoning~\cite{Tran2018Spectral}. SPECTRE strengthens this idea with robust covariance estimation, which amplifies a weak spectral signal when the poisoning ratio is too small for ordinary covariance statistics~\cite{Hayase2021SPECTRE}. Beatrix instead summarizes high-order activation correlations with Gram matrices and learns class-conditional benign statistics, enabling detection of triggers that vary across inputs~\cite{Ma2023Beatrix}. These methods demonstrate that a successful backdoor leaves measurable structure in learned representations, but their supervised formulations commonly exploit labels, predicted classes, or sufficiently coherent within-class feature statistics.

Those conditions are not directly available in an unlabeled contrastive-learning corpus. Global statistics may be dominated by semantic diversity, while class-conditional partitioning based on a compromised encoder can reinforce the association being detected. Passive features can also overlap when the benign and backdoor objectives are tightly coupled~\cite{Li2024Difficulty}. Our approach therefore analyzes paired views of the representation: the original sample describes its learned semantics, while the probe-transformed sample reveals its response to a controlled target association. UMAP preserves local neighborhoods after the training and reference representations are embedded jointly~\cite{McInnes2018UMAP}; DBSCAN then identifies connected high-density regions and noise without pre-specifying the number of clusters~\cite{Ester1996DBSCAN}. Unlike unsupervised anomaly detection that automatically equates the smallest cluster with poisoning, our decision is anchored by the trusted reference center. The dimensionality reduction and density clustering are consequently analysis components, whereas the defense signal originates from controlled semantic probing and trusted reference participation.

\section{Threat Model}\label{sec:threatmodel}
\subsection{Attacker Capabilities}
Let $\mathcal{D}=\mathcal{D}_{b}\cup\mathcal{D}_{p}$ denote the unlabeled pre-training corpus, where $\mathcal{D}_{b}$ and $\mathcal{D}_{p}$ are the benign and poisoned subsets, respectively. The poisoning ratio is
\begin{equation}
    \rho=\frac{|\mathcal{D}_{p}|}{|\mathcal{D}|}.
    \label{eq:poisoningratio}
\end{equation}
The defender does not know $\rho$; we consider the conventional data-poisoning regime in which poisoned samples are a minority. The attacker selects a target semantic class $y_t$, a trigger parameter $\tau$, and a trigger transformation $\mathcal{T}_{\tau}(\cdot)$. The transformation may be a visible patch, a distributed or frequency-domain pattern, or an optimized sample-dependent perturbation. Poisoned examples may contain target-class content or be constructed through an attack-specific optimization, and they participate in pre-training under the same augmentations as benign data.

The attack aims to obtain an encoder $f_{\theta}$ that retains benign representation quality while causing a downstream predictor $g_{\phi}\circ f_{\theta}$ to map trigger-bearing inputs to $y_t$. The attacker may know the contrastive algorithm, encoder architecture, augmentation family, training hyperparameters, diagnostic probe generator, reference semantic, and complete cleansing algorithm, and may optimize poisons through a surrogate training process. In particular, no security claim relies on keeping the probe or reference semantic secret. The attacker can inject or replace samples in $\mathcal{D}$, but cannot alter the defender's cleansing code, the post-cleansing retraining procedure, or the defender-held trusted reference set. The attacker chooses $y_t$ independently of the defender's reference semantic $y_r$; both $y_t\neq y_r$ and the adaptive same-target case $y_t=y_r$ are within scope.

\subsection{Defender Capabilities}
The defender possesses the complete unlabeled corpus $\mathcal{D}$ and can train an inspection encoder, run forward inference, access its representations, apply input transformations, and retrain a final encoder after suspicious samples have been removed. No class labels, poison annotations, trigger pattern, target class, attack identity, or poisoning ratio are available for $\mathcal{D}$. The defense must therefore operate without fixing the number of poisons or semantic clusters in advance.

In addition, the defender holds a trusted benign reference set $\mathcal{R}$ drawn from one manually selected semantic category $y_r$. The reference images may originate from another dataset and are not assumed to match the class composition or acquisition distribution of $\mathcal{D}_{b}$; they need only share a compatible visual modality and input format. We use a reference budget of approximately $|\mathcal{R}|=0.01|\mathcal{D}|$. Labels are used solely to select $y_r$ and never enter a supervised loss or provide poison annotations.

Let $G$ denote the fixed BLTO diagnostic probe associated with the same reference semantic $y_r$. The generator is trained once and reused across datasets and encoder settings. During diagnostic training, the defender applies $G$ to each reference image and trains on $\mathcal{D}\cup G(\mathcal{R})$ with the same stochastic augmentations and unsupervised contrastive objective. Each training and reference sample is traversed once per epoch and receives the same per-sample weighting. After training, the defender compares the original and probe-transformed representations of the training samples and jointly embeds the resulting training and reference information for density-based cleansing. The final encoder is trained from scratch on the retained subset of $\mathcal{D}$; $\mathcal{R}$ is not included in this final retraining stage.

\subsection{Assumptions and Defense Objective}
We make two substantive assumptions. First, $\mathcal{R}$ is benign and protected from poisoning or replacement by the attacker. Second, poisoned samples remain a minority, although their exact prevalence is unknown. We do not assume that the probe generator, the reference semantic, or the cleansing procedure is secret. Probe--attack target alignment is instead treated as an experimental factor: competition between distinct target associations motivates the method, whereas $y_t=y_r$ tests whether the defense remains effective when that competition is weakened or absent.

Let $\delta(x_i)\in\{0,1\}$ be the cleansing decision for $x_i\in\mathcal{D}$, where $1$ denotes a suspicious sample. The estimated poisoned set and the retained corpus are
\begin{equation}
    \widehat{\mathcal{D}}_{p}=\{x_i\in\mathcal{D}:\delta(x_i)=1\},\qquad
    \widehat{\mathcal{D}}_{b}=\mathcal{D}\setminus\widehat{\mathcal{D}}_{p}.
    \label{eq:cleansingsets}
\end{equation}
The immediate objective is to recover as much of $\mathcal{D}_{p}$ as possible while minimizing removal of benign samples. The end-to-end objective is that retraining on $\widehat{\mathcal{D}}_{b}$ substantially reduces attack success on triggered inputs while preserving clean downstream utility. Accordingly, detection precision, recall, and false-positive rate must be reported together with post-cleansing attack success rate and clean accuracy; no single clustering score is treated as sufficient evidence of a successful defense.

\section{Motivation and Key Observations}\label{sec:motivation}
\subsection{Why Passive Representations Are Insufficient}
Most representation-based cleansing methods begin by asking whether a sample is already unusual in the feature space learned from the contaminated corpus. This premise is fragile for contrastive backdoors. A poisoned image still contains ordinary semantic content, and contrastive augmentations encourage its views to remain mutually consistent. Meanwhile, the trigger association is learned as an additional invariant rather than as a replacement for the original semantics. Consequently, benign and poisoned samples may share neighborhoods, losses, and dominant feature directions even when the trained encoder exhibits a high attack success rate~\cite{Li2024Difficulty}. Semantic diversity further complicates global anomaly detection: a rare benign category can be more distant from the feature majority than a poisoned sample whose content belongs to a common category. These considerations explain why the absence of a separate poison cluster in the original representation is not evidence that the corpus is clean.

This limitation motivates a shift from passive inspection to controlled intervention. Instead of asking only where a sample is located, the defender can ask how its representation changes when a known semantic signal is introduced. The intervention supplies a common direction against which samples can be compared, converting an otherwise latent constraint into an observable response. The diagnostic value therefore lies in the difference between the original and intervened views, not in the visual conspicuousness of the probe itself.

\subsection{Observation I: A Semantic Probe Exposes Competing Associations}
BLTO optimizes a trigger that remains associated with a target semantic under contrastive augmentations~\cite{Sun2024BLTO}. We use this property in the opposite role: a pre-trained BLTO generator supplies a controlled reference association rather than an attack against the corpus. For a benign sample, the diagnostic probe introduces a new semantic attraction without competing with a learned backdoor target. For a poisoned sample, the probe-conditioned representation is additionally constrained by the attacker's existing trigger--target association. When the attack and reference targets differ, these two attractions compete, which can reduce or redirect the probe-induced displacement of the poisoned sample relative to benign data.

The resulting mechanism does not require all attacks to share a trigger pattern. Patch, frequency-domain, and optimized triggers differ in appearance and construction, but each successful attack must establish a persistent association in the encoder. The probe is intended to interrogate that learned association rather than reconstruct the original trigger. We therefore use the same diagnostic generator for every training sample and reuse it across datasets and encoders. At the same time, the competition explanation has a clear boundary: if the attack target coincides with the reference semantic, the two signals may reinforce rather than oppose one another. We treat this same-target condition as a required adaptive evaluation rather than excluding it through a secrecy assumption.

\subsection{Observation II: Trusted References Must Participate in Diagnostic Training}
A cross-dataset reference set cannot simply be compared with the contaminated corpus in an unadapted feature space. Acquisition and semantic differences may dominate the comparatively weak backdoor signal, causing a detector to separate data sources rather than poisoning status~\cite{Pan2023ASSET,Chen2025AUTO}. Our design therefore does not reserve $\mathcal{R}$ solely for post-hoc threshold calibration. After applying the diagnostic probe, $G(\mathcal{R})$ participates in the same unsupervised contrastive training as $\mathcal{D}$, under identical augmentations and without a supervised label loss. This exposes the inspection encoder to the reference response while preserving the unlabeled nature of the cleansing problem.

Only one manually selected reference category is used, and its size is approximately $1\%$ of the training corpus. The reference category matches the semantic target of the diagnostic generator, but need not match the attacker's target or the class composition of $\mathcal{D}$. This construction gives the reference points a precise operational role: they identify the region induced by the known diagnostic semantic. Their labels do not identify poisons, and the reference set is removed from the pipeline after cleansing; the final encoder is retrained only on retained training samples.

\subsection{Observation III: Local Topology Supports a Reference-Anchored Decision}
Probe-conditioned differences need not be linearly separable in the original high-dimensional space. We consequently analyze local topology rather than impose a global linear boundary. Training and reference information are embedded jointly into two dimensions with UMAP using cosine distance~\cite{McInnes2018UMAP}; DBSCAN then operates with Euclidean distance in the embedded space~\cite{Ester1996DBSCAN}. Joint fitting is essential because it places the trusted reference points and candidate training samples in one coordinate system instead of comparing independently fitted manifolds.

The reference center is computed as the arithmetic mean of all embedded reference points, including reference points that DBSCAN labels as noise. For each non-noise cluster containing training samples, its center is computed from the training samples in that cluster. The cluster closest to the reference center is interpreted as the benign cluster; all other training clusters and training-sample noise are marked suspicious. If DBSCAN produces only one valid cluster, that cluster is retained and only noise is removed. If no valid cluster exists, cleansing terminates without deleting training data. These rules avoid automatically equating a small cluster with poisoning and make the failure behavior explicit.

Together, these observations yield the working hypothesis evaluated in this study: a reusable semantic probe, a small probed reference set involved in unsupervised diagnostic training, and reference-anchored topology can expose poisoned samples that remain concealed in passive representations. The hypothesis remains subject to three important tests: whether the effect transfers across attack mechanisms and CL algorithms, whether it is stable to reference-set size and manifold randomness, and whether it persists when the attack and diagnostic targets coincide.
\section{Methodology}\label{sec:method}
% TODO
\section{Experimental Evaluation}\label{sec:experiment}
% TODO
\section{Discussion and Limitations}\label{sec:discussion}
% TODO
\section{Conclusion}\label{sec:conclusion}
% TODO

\bibliographystyle{IEEEtran}
\bibliography{references}
\end{document}
